\documentclass[11pt]{article}

\usepackage[margin=1in]{geometry}
\usepackage{amsmath,amssymb}
\usepackage{booktabs}
\usepackage{graphicx}
\usepackage[hidelinks]{hyperref}
\usepackage{natbib}
\usepackage{caption}
\usepackage{enumitem}

\bibliographystyle{plainnat}

\title{\textbf{How Much Does Limit Order Book Microstructure Add\\
Beyond Quote Geometry?}\\[0.4em]
\large Evidence from Cryptocurrency Limit Order Books}

\author{Jude Kriel Ramcharitar\\
\small Independent Researcher\\
\small \texttt{j.k.ramcharitar@gmail.com}}

\date{September 2026 \quad$\cdot$\quad Version 5.0 --- Journal Submission Candidate}

\begin{document}
\maketitle

\begin{abstract}
\noindent
We measure the incremental predictive information in limit order book (LOB)
microstructure after conditioning on three quantities known when a limit order
is placed: quote side, distance from mid, and holding horizon. Using 269{,}330
10-level LOB snapshots collected from Kraken over 70.5 days for BTC/USD,
ETH/USD and SOL/USD, we predict a replay-defined crossing outcome at 300- and
900-second horizons. The benchmark is a 16-cell lookup table containing only
the historical crossing rate for each combination of side, quote offset and
horizon. Under a snapshot-grouped forecasting-ordered split with a 990-second
embargo, gradient-boosted trees using 63 microstructure features improve AUC by
0.0479 on BTC/USD and 0.0596 on ETH/USD, while SOL/USD is 0.0105 below the
lookup reference. The SOL result is sensitive to temporal distribution shift:
under an exploratory blocked diagnostic that increases train-test price-range
overlap, the microstructure gap is positive for all three assets, and remains
positive in all 30 asset-by-phase evaluations across ten block rotations. These
blocked results are interpreted as evidence of information content rather than
prospective forecasting performance because training observations may post-date
test observations. Model-class comparisons are representation-dependent and do
not support a general ranking between gradient-boosted trees and an LSTM. The
results show that LOB state contains measurable incremental information beyond
quote geometry in this sample, while also illustrating how split design and
pipeline validation can materially change conclusions.

\vspace{0.8em}
\noindent\textbf{Keywords:} limit order book, market microstructure, crossing
probability, gradient boosting, predictive information, reproducibility

\vspace{0.3em}
\noindent\textbf{JEL:} G12, G14, C45, C61
\end{abstract}

\vspace{1em}

\section{Introduction}

A trader placing a limit order chooses the side, the distance from mid, and the
length of time the order is allowed to rest. Those choices mechanically affect
whether the quoted price is subsequently reached. A useful evaluation of LOB
predictive models should therefore distinguish information that comes from the
state of the book from information already contained in quote geometry.

This paper asks a deliberately narrow question: conditional on side, quote
offset and horizon, how much additional information about a future price
crossing is contained in observable LOB microstructure?

The distinction matters because a model can achieve a respectable headline AUC
without adding much beyond variables that were already chosen by the trader.
Existing work demonstrates substantial predictive content in LOB data and
execution-related models \citep{zhang2019,maglaras2022,arroyo2024}, but the
incremental value of book state is harder to interpret when the reference model
already embeds placement choices. We therefore benchmark against an explicit
lookup reference that uses only side, quote offset and horizon.

Our primary evaluation is forecasting-ordered. On BTC/USD and ETH/USD,
gradient-boosted trees using 63 microstructure features improve AUC over the
lookup reference by 0.0479 and 0.0596. On SOL/USD the corresponding gap is
-0.0105. The negative SOL result coincides with substantial temporal
distribution shift: train and test price ranges overlap by only 0.251, compared
with 0.129 for BTC and 0.073 for ETH. Because absolute price can become a proxy
for partition membership over a trending sample, we separately use a blocked
split as an information-content diagnostic. That split raises price-range
overlap and produces positive microstructure gaps on all three assets. Across
ten block-phase rotations, all 30 asset-by-phase gaps remain positive, with
means of 0.0473, 0.0382 and 0.0346 for BTC, ETH and SOL.

The blocked analysis is intentionally secondary. It is not a prospective
forecasting test because training observations may occur after test
observations. Its role is to diagnose whether the temporal result is consistent
with absent microstructure signal or with sensitivity to regime composition and
distribution shift. The temporal split remains the primary test for any
forecasting interpretation.

A second contribution is methodological. Several implementation defects during
development produced plausible-looking results, including errors in horizon
resolution, timestamp units, model conditioning and target specification. The
final pipeline includes simple references, positive controls and assertions on
split properties. We report these validation checks because the defects were
not identifiable from headline metrics alone.

The paper makes four contributions:
\begin{enumerate}[itemsep=2pt,topsep=4pt]
\item It directly measures the incremental predictive value of LOB
microstructure over quote geometry on three cryptocurrency pairs.
\item It separates a forecasting-ordered evaluation from an exploratory blocked
information-content diagnostic and quantifies sensitivity across ten block
phases.
\item It compares gradient-boosted trees and an LSTM while holding feature
representation fixed, finding no general model-class ordering.
\item It documents validation checks that detected plausible but incorrect
results in the research pipeline.
\end{enumerate}

The claims are intentionally limited. The outcome is a replay-defined crossing,
not a realised fill; the paper does not model queue position, transaction costs,
adverse selection or profitability; AUC measures ranking rather than
calibration; and the evidence comes from one venue, one 70.5-day window and
three cryptocurrency pairs.

\section{Related Work}

\subsection{Execution and crossing prediction}

\citet{lo2002} established econometric foundations for execution probability
modelling. \citet{maglaras2022} model execution as a time-to-fill distribution,
and \citet{arroyo2024} use a convolutional-transformer survival framework for
limit-order execution. Those studies use substantially finer market data than
ours. Our snapshot feed has median gaps of 11.5 to 48.0 seconds, so our crossing
outcome should not be interpreted as a substitute for tick-level execution
modelling. The contribution here is instead the explicit separation between the
mechanical information in quote placement and the incremental information in
book state.

\subsection{Financial machine-learning baselines and validation}

The need for simple references and careful temporal validation is well known in
financial machine learning. \citet{lopezdeprado2018} discusses leakage from
serial dependence and motivates purging and embargoing. We adopt a
snapshot-grouped split and a 990-second embargo because labels look as far as
900 seconds into the future. Our emphasis is not that trivial baselines are
novel, but that they provide a direct interpretation of what the microstructure
model adds over placement geometry.

\subsection{Model class on tabular financial features}

Gradient-boosted trees remain competitive with neural networks on many tabular
problems \citep{grinsztajn2022,shwartz2022}. Our engineered per-snapshot
features are tabular in structure, while the LSTM additionally receives a
50-snapshot history. We therefore treat the model comparison as descriptive and
hold feature representation fixed when comparing classes.

\section{Data}
\label{sec:data}

\subsection{Collection}

LOB data were collected from Kraken WebSocket API v2 using the depth-10
\texttt{book} channel for BTC/USD, ETH/USD and SOL/USD. Updates were timestamped
at receipt and stored with the full 10-level ladder in Parquet. Collection ran
from 15 June to 24 August 2026, a span of 70.5 days. A 14-day collector outage
occurred during the window, so mean inter-update gaps include downtime; median
gaps better describe normal feed cadence.

\begin{table}[ht]
\centering
\caption{Collected data. Label rows reflect observable replay-defined crossing
outcomes across four offsets, two horizons and two sides.}
\label{tab:data}
\begin{tabular}{lrrrrr}
\toprule
Asset & Snapshots & Tick & Median gap & Mean gap$^{*}$ & Label rows \\
\midrule
BTC/USD & 111{,}002 & 0.10 & 11.5\,s & 54.9\,s & 1{,}760{,}728 \\
ETH/USD & 34{,}380 & 0.01 & 48.0\,s & 177.2\,s & 515{,}576 \\
SOL/USD & 123{,}948 & 0.01 & 26.2\,s & 49.1\,s & 1{,}971{,}128 \\
\midrule
Total & 269{,}330 & & & & 4{,}247{,}432 \\
\bottomrule
\end{tabular}

\vspace{0.4em}
{\footnotesize $^{*}$Includes the collector outage.}
\end{table}

\subsection{Outcome construction}
\label{sec:labels}

For each snapshot we simulate limit orders at offsets of 1, 2, 5 and 10 basis
points from mid on each side, rounded to the venue tick, and record whether the
opposite best quote reaches or crosses that price within 300 or 900 seconds. A
buy at price $p$ is recorded as crossed if
$\min(\text{best ask}) \leq p$ within the horizon; a sell is crossed if
$\max(\text{best bid}) \geq p$.

We call this a \emph{replay-defined crossing}. It uses subsequent book
snapshots, not trade or order-level execution data, so it does not establish
that a resting order actually filled. Rows for which the horizon is not
observable are dropped rather than labelled as non-crossings. Observability is
88.7\% at worst (ETH/USD at 300 seconds) and 98.8\% or above at 900 seconds.

Offsets are expressed in basis points rather than ticks because venue ticks are
not economically comparable across assets. Horizons are resolved against real
timestamps rather than snapshot counts.

\begin{table}[ht]
\centering
\caption{Crossing rate by horizon and quote offset. The strong monotone
structure motivates conditioning the reference on quote geometry.}
\label{tab:marginals}
\begin{tabular}{llrrrr}
\toprule
Asset & Horizon & 1\,bps & 2\,bps & 5\,bps & 10\,bps \\
\midrule
BTC/USD & 300\,s & 0.697 & 0.622 & 0.433 & 0.237 \\
& 900\,s & 0.815 & 0.766 & 0.626 & 0.440 \\
ETH/USD & 300\,s & 0.603 & 0.540 & 0.398 & 0.245 \\
& 900\,s & 0.740 & 0.696 & 0.582 & 0.432 \\
SOL/USD & 300\,s & 0.577 & 0.517 & 0.387 & 0.242 \\
& 900\,s & 0.753 & 0.712 & 0.612 & 0.466 \\
\bottomrule
\end{tabular}
\end{table}

\subsection{Features}

We construct 109 features from each snapshot, including raw and normalised price
and volume levels, mid-price, spread, spread in basis points, order-book and
depth imbalances, rolling statistics and lagged differences.

The headline \emph{microstructure} subset contains 63 features matching spread,
imbalance, order-book imbalance, depth, volume or volatility terms while
excluding variables whose names contain \texttt{price}. It also excludes five
lagged mid-price return features. The subset therefore removes direct price
levels and directional lagged returns, although spreads and volatility remain
functions of prices. A 41-feature price-named subset is used only for
diagnostics. The full representation contains all 109 features.

\section{Method}
\label{sec:method}

\subsection{Reference model}

The benchmark is a 16-cell lookup table: the mean training-set crossing rate for
each (offset, horizon, side) combination, applied unchanged to the test set. It
contains no LOB state and has no learned functional form. Its purpose is to
capture information mechanically available from quote geometry before asking
whether the book adds anything further.

\subsection{Models}

\textbf{Gradient-boosted trees.} XGBoost with 400 estimators, depth 6, learning
rate 0.05 and subsample 0.8, using the current snapshot plus the three
conditioning variables.

\textbf{Binary LSTM.} A two-layer LSTM with hidden dimension 64 over a
50-snapshot window. The conditioning vector is concatenated before a single
logit head trained with binary cross-entropy. The LSTM receives a temporal
window while the tree receives only the current snapshot, so the comparison is
not interpreted as a controlled architecture tournament.

\subsection{Sampling, splitting and embargo}
\label{sec:split}

For computational comparability, each asset uses a uniform random sample without
replacement of 200{,}000 label rows where available, with a fixed sampling seed,
after which temporal order is restored. The same sampled rows are used across
model arms.

Each snapshot generates up to 16 label rows with identical snapshot features.
Splits are therefore grouped by snapshot to prevent one snapshot from appearing
in multiple partitions. Because labels look up to 900 seconds forward, every
partition boundary carries a 990-second embargo.

The \textbf{temporal split}, our primary forecasting evaluation, assigns the
first 70\% of snapshots to training, the next 15\% to validation and the final
15\% to test.

The \textbf{blocked split}, used only as a secondary information-content
diagnostic, divides the timeline into 150 contiguous blocks assigned in a
repeating train/validation/test pattern. This increases overlap in market
regimes across partitions but permits training observations to post-date test
observations. It therefore cannot be interpreted as a prospective trading or
forecasting evaluation.

\subsection{Metric}

We report ROC AUC for the binary crossing outcome, computed with ties credited
0.5. AUC measures ranking rather than calibration and does not imply economic
value.

\section{Results}
\label{sec:results}

\subsection{Primary forecasting-ordered results}

\begin{table}[ht]
\centering
\caption{Crossing AUC under the temporal split, three model seeds. Seed
standard deviations measure model-initialisation variation only.}
\label{tab:main}
\small
\begin{tabular}{lrrrrr}
\toprule
Asset & Lookup & GBT micro (63) & GBT all (109) & LSTM micro (63) & LSTM all (109) \\
\midrule
BTC/USD & 0.6828 & 0.7307 & 0.7262 & 0.6857 & 0.6903 \\
& & \tiny$\pm$0.0006 & \tiny$\pm$0.0012 & \tiny$\pm$0.0021 & \tiny$\pm$0.0043 \\
& & $+0.0479$ & $+0.0434$ & $+0.0029$ & $+0.0075$ \\
\addlinespace
ETH/USD & 0.6597 & 0.7193 & 0.6908 & 0.6983 & 0.6669 \\
& & \tiny$\pm$0.0013 & \tiny$\pm$0.0010 & \tiny$\pm$0.0046 & \tiny$\pm$0.0036 \\
& & $+0.0596$ & $+0.0311$ & $+0.0386$ & $+0.0072$ \\
\addlinespace
SOL/USD & 0.6673 & 0.6568 & 0.7056 & 0.7043 & 0.7043 \\
& & \tiny$\pm$0.0035 & \tiny$\pm$0.0010 & \tiny$\pm$0.0046 & \tiny$\pm$0.0074 \\
& & $-0.0105$ & $+0.0383$ & $+0.0370$ & $+0.0370$ \\
\bottomrule
\end{tabular}
\end{table}

The headline microstructure result is positive on BTC/USD and ETH/USD: the GBT
using the 63-feature subset gains 0.0479 and 0.0596 AUC over the lookup
reference. SOL/USD is the exception at -0.0105. Adding the non-microstructure
columns lowers GBT performance on BTC and ETH but raises it on SOL, suggesting
that the SOL sample behaves differently under the forecasting-ordered split.
Section~\ref{sec:regime} examines this sensitivity.

Seed standard deviations are small relative to the BTC and ETH gaps, but they
measure only model-initialisation variation. They are not confidence intervals
for the data-generating process, and the label rows are not treated as
independent observations for inferential purposes because multiple rows arise
from each snapshot and adjacent snapshots are serially dependent.

\subsection{Model-class comparison}
\label{sec:modelclass}

Comparisons between GBT and LSTM are made within the same feature
representation rather than by selecting each model's best test-set arm.

\begin{table}[ht]
\centering
\caption{GBT minus LSTM AUC under the temporal split, holding feature
representation fixed. Positive values favour GBT.}
\label{tab:modelclass}
\begin{tabular}{lrrrr}
\toprule
& \multicolumn{2}{c}{Microstructure (63)} & \multicolumn{2}{c}{All (109)} \\
\cmidrule(lr){2-3}\cmidrule(lr){4-5}
Asset & GBT $-$ LSTM & Leader & GBT $-$ LSTM & Leader \\
\midrule
BTC/USD & $+0.0450$ & GBT & $+0.0359$ & GBT \\
ETH/USD & $+0.0210$ & GBT & $+0.0239$ & GBT \\
SOL/USD & $-0.0475$ & LSTM & $+0.0013$ & GBT \\
\bottomrule
\end{tabular}
\end{table}

The ranking is representation-dependent. Trees lead on BTC and ETH under both
representations. On SOL, the LSTM leads by 0.0475 on the microstructure subset,
while GBT leads by only 0.0013 on the full representation, less than the
reported seed dispersion of the SOL LSTM arm. We therefore draw no general
model-class conclusion from these data. The ETH LSTM-all arm is also likely
undertrained: its training loss moved only 0.0078, versus 0.1593 for LSTM-micro
on the same data.

\section{Split-Scheme Sensitivity}
\label{sec:regime}

\subsection{Temporal distribution shift}

The negative SOL microstructure gap motivated a diagnostic of train-test price
range overlap. Under the temporal split, overlap is low on all three assets.

\begin{table}[ht]
\centering
\caption{Overlap between training and test price ranges. Low overlap means
absolute price can become informative about partition membership.}
\label{tab:overlap}
\begin{tabular}{lrr}
\toprule
Asset & Temporal & Blocked \\
\midrule
BTC/USD & 0.129 & 0.815 \\
ETH/USD & 0.073 & 0.855 \\
SOL/USD & 0.251 & 0.745 \\
\bottomrule
\end{tabular}
\end{table}

The blocked split substantially increases overlap, but changes more than price
range alone: it also changes regime composition, temporal distance between
partitions and the direction of information flow. We therefore do not interpret
it as isolating a price-level mechanism.

\subsection{Blocked information-content diagnostic}

\begin{table}[ht]
\centering
\caption{Gap over lookup by feature group and split scheme, three seeds. The
blocked split is diagnostic rather than prospective.}
\label{tab:regime}
\begin{tabular}{lrrrrrr}
\toprule
& \multicolumn{2}{c}{Microstructure (63)} & \multicolumn{2}{c}{Price-named (41)}
& \multicolumn{2}{c}{All (109)} \\
\cmidrule(lr){2-3}\cmidrule(lr){4-5}\cmidrule(lr){6-7}
Asset & Temporal & Blocked & Temporal & Blocked & Temporal & Blocked \\
\midrule
BTC/USD & $+0.0479$ & $+0.0479$ & $+0.0384$ & $+0.0383$ & $+0.0434$ & $+0.0456$ \\
ETH/USD & $+0.0596$ & $+0.0491$ & $-0.0001$ & $+0.0370$ & $+0.0312$ & $+0.0513$ \\
SOL/USD & $-0.0105$ & $+0.0274$ & $+0.0312$ & $+0.0115$ & $+0.0383$ & $+0.0295$ \\
\bottomrule
\end{tabular}
\end{table}

BTC and ETH remain positive under both schemes. SOL changes from -0.0105 under
the temporal split to +0.0274 under blocking. We interpret that reversal as
split sensitivity under temporal distribution shift, not as proof of a specific
causal mechanism. Price-named features also behave inconsistently across assets,
so headline claims rely on the 63-feature microstructure subset.

\subsection{Block-phase rotations}
\label{sec:rotation}

The blocked assignment has an arbitrary phase. We therefore rotate the phase
through all ten assignments while holding the 150 blocks, embargo, feature
representation and model specification fixed. Test sets are disjoint across the
ten phases for each asset, although training sets overlap heavily and the
rotations are not statistically independent.

\begin{table}[ht]
\centering
\caption{Microstructure gap over lookup across ten blocked phases.}
\label{tab:rotation}
\begin{tabular}{lrrrrr}
\toprule
Asset & Mean & SD & Min & Max & Positive \\
\midrule
BTC/USD & $+0.0473$ & 0.0052 & $+0.0378$ & $+0.0529$ & 10/10 \\
ETH/USD & $+0.0382$ & 0.0144 & $+0.0130$ & $+0.0550$ & 10/10 \\
SOL/USD & $+0.0346$ & 0.0103 & $+0.0133$ & $+0.0463$ & 10/10 \\
\bottomrule
\end{tabular}
\end{table}

All 30 asset-by-phase evaluations are positive. The across-phase dispersion is
substantially larger than model-seed dispersion, which underscores that
partition choice matters more than stochastic model initialisation in this
experiment. These standard deviations are descriptive sensitivity measures,
not independent-sample standard errors.

Price-range overlap itself does not show a consistent association with the gap
across the ten phases: correlations are -0.17, +0.33 and -0.14 for BTC, ETH and
SOL. This weak and inconsistent pattern is another reason not to attribute the
SOL reversal to price overlap alone.

\section{Validation, Failure Modes, and Reproducibility Checks}
\label{sec:defects}

During development, several pipeline failures produced outputs that would have
looked plausible without targeted diagnostics. We report them because they
motivate the final validation design.

\textbf{Conditioning omission.} Quote offset was initially absent from the
conditioning vector. Label rows from the same snapshot therefore shared model
inputs while carrying different targets. A lookup-table positive control exposed
the failure because the reference could outperform the learned model.

\textbf{Incorrect horizon resolution.} A synthetic 100-millisecond interval was
initially applied to real, irregularly sampled data. Nominal millisecond
horizons then corresponded to tens or hundreds of seconds in calendar time. The
final pipeline resolves horizons against timestamps.

\textbf{Timestamp-unit error.} Microsecond timestamps were once interpreted as
nanoseconds, shrinking measured intervals by a factor of 1000. A median-gap
assertion now fails when the recovered cadence is inconsistent with the source
feed.

\textbf{Target-head mismatch.} Binary labels were initially passed to a
survival-style decoder with a constant observed time, producing a fixed loss and
AUC of 0.5000. A planted-signal positive control verified that the corrected
binary head can learn when signal is present.

\textbf{Split-proportion error.} Partition boundaries based on elapsed time
rather than snapshot count produced realised proportions far from the intended
70/15/15 split. The final pipeline asserts realised grouped proportions.

The common lesson is that plausible headline metrics are not sufficient
validation. The final workflow therefore runs the trivial reference before the
main model, includes a planted-signal positive control, verifies timestamp
cadence and asserts split properties and embargo conditions.

\section{Limitations}

\textbf{Crossing is not execution.} The outcome is reconstructed from future
book snapshots. A price crossing does not establish that a resting order filled,
especially for an order arriving late in the queue.

\textbf{Transient crossings can be missed.} With median inter-update gaps of
11.5 to 48.0 seconds, a crossing can occur and reverse between observations.
The direction of any resulting bias is not established.

\textbf{Observability is state-dependent.} Labels require at least one future
book observation within the horizon. Quiet states may therefore be
underrepresented. In the worst case, ETH/USD at 300 seconds, 11.3\% of
potential observations are dropped.

\textbf{No economic evaluation.} We do not model fees, spread capture, queue
position, adverse selection or trading profitability. AUC is not a measure of
expected return.

\textbf{No calibration analysis.} The paper evaluates ranking only.

\textbf{Limited external validity.} The sample covers one venue, three assets,
one 70.5-day period and two horizons.

\textbf{Limited model search.} The deep-learning comparison uses one LSTM
architecture and the ETH full-feature LSTM arm appears undertrained.

\textbf{Uncertainty is only partially characterised.} Model-seed dispersion and
blocked-phase dispersion are reported, but neither is a general sampling
distribution for the forecasting-ordered effect. The blocked phases share
training data and the temporal evaluation admits only one chronological test
period in this sample.

\section{Conclusion}

After conditioning on quote side, distance from mid and holding horizon, LOB
microstructure contains measurable incremental predictive information in this
sample. Under the primary forecasting-ordered evaluation, the 63-feature
microstructure model improves AUC by 0.0479 on BTC/USD and 0.0596 on ETH/USD,
while SOL/USD is 0.0105 below the lookup reference. The SOL result is sensitive
to partition design: under a blocked information-content diagnostic the gap is
positive, and it remains positive across all ten block phases.

The evidence therefore supports a narrower conclusion than a generic claim that
``LOB models predict fills.'' Book state adds information beyond quote geometry
on BTC and ETH in prospective order, while the SOL result shows that the sign
and magnitude of the measured increment can depend on temporal distribution
shift. The blocked experiment indicates that microstructure information is
present in SOL as well, but it does not establish that the information is
prospectively exploitable.

The practical next question is economic rather than statistical: whether the
incremental ranking information has value after queue position, fees, spread,
adverse selection and execution uncertainty. This paper does not answer that
question.

\section*{Data and Code Availability}
\addcontentsline{toc}{section}{Data and Code Availability}

The collection pipeline, label builder, model code, verification scripts and
reported diagnostics are available at
\url{https://github.com/juderamcharitar/latency-compensating-mm}.

\section*{Disclosure of AI Assistance}
\addcontentsline{toc}{section}{Disclosure of AI Assistance}

The author used Claude (Anthropic) extensively for implementation, diagnostic
analysis, drafting and revision. AI-generated code contributed to several of the
pipeline failures described above and AI-assisted diagnosis contributed to
identifying several of them. The author designed the study, directed the
analysis, made the substantive research decisions, verified the reported
results, and accepts responsibility for the content and any errors.

\section*{Acknowledgements}
\addcontentsline{toc}{section}{Acknowledgements}

This work was conducted independently, without institutional funding or
supervision. The author declares no conflicts of interest and holds no positions
in any asset discussed.

\bibliography{../paper_v4/refs}

\end{document}
